% page 516 of the facsimile (image p-515 of the scan)
% block 157.5 x 242.0 mm ; left margin 8.0 mm
\pgc{-12.700mm}{0.000mm}{
\l{\cel{3}508.}
\l{}
\l{\sou{saja}.\cel{2}Qua judikas yuste pri la kozi, dicernas lo vera de}
\l{\cel{3}lo falsa, lo exakta de lo semblanta (segun ke ta fakul-}
\l{\cel{3}tato esas komuna ye omna homi), agas segun raciono e mo-}
\l{\cel{3}dere - e di qua la konduto esas delikate regulkonforma.}
\l{\cel{3}- EFIS.}
\l{}
\l{\sou{sako}. I. Peco de ledro, de stofo faldita mediane ; klozita}
\l{\cel{3}(per suturo o per glutino) laterale ; kun aperturo supre}
\l{\cel{3}por posibligar la enpozo di objekti.II. (\sou{fiziol.}) Kavajo}
\l{\cel{3}en la korpo, cirkondata da parieto membrana. - DEFIS.}
\l{}
\l{\sou{sakarimetro}.\cel{1}\sou{(kemio}).Aparato qua mezuras, per la deviaco di}
\l{\cel{3}la radio polarigita, la quanto de sukro quan kontenas li-}
\l{\cel{3}quido. - DEFIS.}
\l{}
\l{\sou{sakarino}. (\sou{kemio).}\cel{2}Laktono fuzebla ye 100 gradi C., qua sa-}
\l{\cel{3}poras bitre, obtenita per la efiko di kalko sur glukoso}
\l{\cel{3}o levuloso. -\cel{1}\sou{Simbolo kemiala}\cel{1}: C6 H4 / CO SO2/ MH.- DEFIS.}
\l{}
\l{\sou{sakra}. (\sou{religio}).I. Di qua la karaktero esas neviolacenda}
\l{\cel{3}pro ke konsakrita a la oficio deala. - II. Egardata kom}
\l{\cel{3}neviolacenda. - EFIS.}
\l{}
\l{\sou{sakramento}. (\sou{religio katolika}).Singla de la agi, institucita}
\l{\cel{3}da Iesu Kristo, por la santigo di la anmi, signo sentebla}
\l{\cel{3}di efekto spiritala, destinita naskigar vivo de gracio}
\l{\cel{3}(bapto, penitenco) od augmentar lu (konfirmo, eukaristio,}
\l{\cel{3}mariajo, ordo, unciono extrema). - DEFIRS.}
\l{}
\l{\sou{sakrifikar}. I. (\sou{trans.}) (\sou{religio}).1. Kombustar (kozo quan}
\l{\cel{3}onu egardas kom tre precoza) sur altaro konsakrita a deo,}
\l{\cel{3}por honorizar o flexar lu.- 2. Ofrar vivanto quan onu imo-}
\l{\cel{3}las por la deo. - II. Perdar volunte persono o kozo quan}
\l{\cel{3}onu amas o prizas rispektive, po (ulu, ulo). - EFIS.}
\l{}
\l{\sou{sakrilejar}. (\sou{trans.)}\cel{1}(\sou{religio}).Violacar ulo sakra. - DEFIS.}
\l{}
\l{\sou{sakristo}. (\sou{katolikismo}).La persono qua su okupas ye ta parto}
\l{\cel{3}acesora di kirko ube esas depozita la vazi e la ornivi}
\l{\cel{3}sakra, ed ube la oficionti metas sua vesti sacerdotala ;}
\l{\cel{3}e di ta parto di kirko, ube la jus-mariajiti e lia testi}
\l{\cel{3}venas signatar sur la registri di la parokio, e gratulesar}
\l{\cel{3}da la amiki qui asistis la mariajo. - DEFIS.}
\l{}
\l{\sou{sakrumo}. (\sou{anat.})\cel{2}Osto di la pelvo, quan la antiqui ofris}
\l{\cel{3}sakrifike. - EFIS.}
\l{}
\l{\sou{salo}. I. Na Cl, substanco friabla, dissolvebla en aquo, qua}
\l{\cel{3}saporas pikante, uzata kom kondimento, por konservar la}
\l{\cel{3}nutrivi, por impedar lia putro. - II. (\sou{kemio}).Kombinajo}
\l{\cel{3}de acido kun oxido. - DEFIRS.}
\l{}
\l{\sou{salado}. Disho ek herbi o legumi kun (kom spici o kondimenti)}
\l{\cel{3}salo, pipro, oleo e vinagro. - DEFIRS.}
}
